\chapter{数字经济中的计量与因果推断方法}

\begin{examguide}
\textbf{本章考情}：本章为方法章（\important{专硕·扩展}），人大等考计量经济学的院校要求掌握 OLS、DID、IV、RDD 的推导逻辑，本章亦为学硕计量经济学与复试论文讨论的基础（\important{学硕·热点}）。数字经济实证研究的识别策略（DID、IV、RDD）是复试问答与论文写作的高频内容。

\textbf{核心考点}：（1）OLS 估计量与高斯-马尔可夫假设；（2）R² 与 t 检验的逻辑；（3）内生性三来源与遗漏变量偏误公式；（4）DID 估计量与平行趋势假设；（5）事件研究图与平行趋势检验；（6）IV/2SLS、LATE 与弱工具变量；（7）RDD 的断点跳跃与局部效应；（8）PSM 与合成控制的分工；（9）实证案例的四步流程（分析问题→建模→求解→分析）；（10）熵值法指数构造；（11）平台 A/B 实验的随机化逻辑。

\textbf{本章主线}：回归基础（13.1）→ 识别策略（13.2）→ 案例教学（13.3）。
\end{examguide}

平台行为是否损害竞争、隐私规制是否抑制创新、数字服务税是否改变价格结构，第 12 章末尾提出的这些命题都需要计量工具来检验，也都归结为两个基本问题：实证结论从哪里来，什么时候可信。本章依次展开：13.1 节建立回归分析的基础语言，13.2 节介绍主流的因果识别策略，13.3 节回到数字经济本身，用实证案例演示从分析问题到结论的完整研究流程。方法与结论不可分：实证结论的可信度，取决于支撑它的识别策略是否成立。

\section{计量基础回顾}

因果推断的讨论需要一套共同语言。本章的方法对象是数字经济实证研究：用数据检验数字化的效应（生产率、收入、竞争），而一切实证结论都始于回归。本节回顾两个基础概念：OLS 估计量回答"结论从哪里来"，内生性回答"结论什么时候不可信"。

\subsection{OLS 估计量}

OLS 是数字经济实证研究的最小单位：估计"数字化程度对绩效的影响"这类线性关系，标准工具就是 OLS。理解 OLS 的推导，才能理解后续所有识别策略为什么是对 OLS 的"修正"。

\begin{definition}[普通最小二乘法（Ordinary Least Squares, OLS）]
普通最小二乘法是以残差平方和最小化为准则估计线性模型系数的方法：在一元模型 $y_i = a + b x_i + \varepsilon_i$ 中选取截距与斜率 $\hat a$、$\hat b$，使 $\sum_{i=1}^{n}(y_i - a - b x_i)^2$ 达到最小。
\end{definition}

\begin{derivation}{OLS 的正规方程}{}
\textbf{OLS 的正规方程推导。}设一元线性回归模型
\[
y_i = a + b x_i + \varepsilon_i, \qquad i = 1, \dots, n,
\]
其中 $y_i$ 为被解释变量，$x_i$ 为解释变量，$\varepsilon_i$ 为扰动项，待估参数为截距 $a$ 与斜率 $b$。OLS 的思路分两步：先构造一个衡量拟合偏离的准则函数，再找使该函数最小的系数。准则函数取残差平方和
\[
SSR(a,b) = \sum_{i=1}^{n} (y_i - a - b x_i)^2,
\]
它度量给定系数 $a$、$b$ 时拟合直线 $a + b x_i$ 与真实观测 $y_i$ 的总偏离，取平方是为了让正负残差同向累积、且方便求导。\textbf{第一步，一阶条件。}最小化问题的解满足偏导数为零：对 $a$ 求偏导并令其为零，
\[
\frac{\partial SSR}{\partial a} = -2 \sum_{i=1}^{n} (y_i - \hat a - \hat b x_i) = 0,
\]
对 $b$ 求偏导并令其为零，
\[
\frac{\partial SSR}{\partial b} = -2 \sum_{i=1}^{n} x_i (y_i - \hat a - \hat b x_i) = 0.
\]
这两个一阶条件合称\textbf{正规方程}。它们的含义：记残差 $\hat\varepsilon_i = y_i - \hat a - \hat b x_i$，第一个条件说残差之和为零，第二个条件说残差与解释变量的样本协方差为零，拟合已把 $y$ 中所有与 $x$ 相关联的成分提取干净，残差中不再剩任何可被 $x$ 解释的信息。\textbf{第二步，求解。}由第一个条件解出截距 $\hat a = \bar y - \hat b \bar x$，代回第二个条件并利用和式性质 $\sum_i (x_i - \bar x) = 0$ 化简，得到斜率
\[
\hat b = \frac{\sum_i (x_i - \bar x)(y_i - \bar y)}{\sum_i (x_i - \bar x)^2} = \frac{\widehat{Cov}(x,y)}{\widehat{Var}(x)},
\]
即斜率等于 $x$ 与 $y$ 的样本协方差除以 $x$ 的样本方差：分子是二者共同变动的部分，分母是 $x$ 自身的变异，两者之比给出"$x$ 每变动一单位、$y$ 平均变动多少"。截距 $\hat a = \bar y - \hat b \bar x$ 保证拟合直线穿过样本均值点 $(\bar x, \bar y)$。\textbf{第三步，无偏性。}将真实模型 $y_i = a + b x_i + \varepsilon_i$ 代入斜率公式，把扰动项分离出来：
\[
\hat b = b + \frac{\widehat{Cov}(x, \varepsilon)}{\widehat{Var}(x)},
\]
估计量等于真实参数加一个噪声项，无偏性要求噪声项的期望为零。两端取期望，\important{无偏性的关键假设为 $\mathbb{E}[\varepsilon_i \mid x_i] = 0$（零条件均值）}：扰动项的条件均值不随解释变量取值变化。此假设成立时 $\mathbb{E}[\hat b] = b$，而所有内生性问题都是这一假设的违背。多元情形的矩阵形式为 $\hat\beta = (X'X)^{-1} X'y$，其中 $X$ 以各解释变量为列、$y$ 为被解释变量向量，考计量经济学的院校要求掌握，其余读者只需知道这是本框两个正规方程的推广。
\end{derivation}

推导之外，OLS 有两层直观含义。\textbf{拟合含义}：正规方程要求残差之和为零、残差与 $x$ 的样本协方差为零，拟合直线因此穿过样本中心，并把 $y$ 中随 $x$ 变动的部分全部拿走，剩下的是与 $x$ 无关的噪声。\textbf{矩条件含义}：零条件均值假设的推论是 $\mathbb{E}[x_i \varepsilon_i] = 0$，正规方程正是这一总体条件的样本版本，这一视角是理解 IV（13.2.3 节）的桥梁。第 1 章验证摩尔定律的半对数回归，就是 OLS 最简单的应用：以 $\ln N$ 对 $t$ 回归检验指数趋势。

线性模型的有效性问题由高斯-马尔可夫定理回答。OLS 的完整假设集共五条：\textbf{线性设定}（$y$ 是参数的线性函数，如 $y = a + bx + \varepsilon$）、\textbf{随机抽样}（样本独立同分布）、\textbf{无完全共线}（一元情形要求 $x$ 有正的样本变异，多元情形要求任一解释变量都不能被其余解释变量完全线性表示）、\textbf{零条件均值}（$\mathbb{E}[\varepsilon_i \mid x_i] = 0$）、\textbf{同方差}（$Var(\varepsilon_i \mid x_i) = \sigma^2$ 不随 $i$ 变化）。五条同时成立时，OLS 是最优线性无偏估计量（Best Linear Unbiased Estimator, BLUE），即高斯-马尔可夫定理。五条假设是分层的：\important{无偏性只需要前四条（含零条件均值），BLUE 需要全部五条}；若零条件均值不成立，OLS 无论样本多大都不收敛到真实参数，这正是 13.1.2 节内生性问题的要害。

拟合优度与显著性检验是回归输出的两个基本量。\textbf{拟合优度}回答"$x$ 在多大程度上解释了 $y$ 的变异"，由 $R^2$ 度量；\textbf{显著性检验}回答"单个系数的估计是否在统计上区别于零"，由 $t$ 统计量度量。

\begin{derivation}{R² 与 t 检验}{}
\textbf{$R^2$ 的平方和分解。}先定义三个平方和：总平方和 $TSS$ 度量 $y$ 的总变异，回归平方和 $ESS$ 度量拟合值 $\hat y$ 解释掉的变异，残差平方和 $RSS$ 度量未被解释的变异：
\[
TSS = \sum_{i=1}^{n}(y_i - \bar y)^2, \qquad
ESS = \sum_{i=1}^{n}(\hat y_i - \bar y)^2, \qquad
RSS = \sum_{i=1}^{n} \hat\varepsilon_i^2.
\]
三者存在精确的分解关系，其推导依赖含截距模型带来的两条正交性。其一，13.1.1 节正规方程的第一个条件即 $\sum_i \hat\varepsilon_i = 0$，残差和为零；其二，由 $\hat y_i = \hat a + \hat b x_i$ 知 $\sum_i \hat y_i \hat\varepsilon_i = \hat a \sum_i \hat\varepsilon_i + \hat b \sum_i x_i \hat\varepsilon_i = 0$，两项分别由正规方程的两个条件消去，拟合值与残差正交。把 $y_i - \bar y = (\hat y_i - \bar y) + \hat\varepsilon_i$ 代入 $TSS$ 并展开：
\[
TSS = \sum_i [(\hat y_i - \bar y) + \hat\varepsilon_i]^2 = ESS + RSS + 2\sum_i (\hat y_i - \bar y)\hat\varepsilon_i,
\]
交叉项 $2\sum_i (\hat y_i - \bar y)\hat\varepsilon_i = 2\sum_i \hat y_i \hat\varepsilon_i - 2\bar y \sum_i \hat\varepsilon_i = 0$，两条正交性恰好把交叉项消掉，得到平方和分解：
\[
TSS = ESS + RSS.
\]
由此定义拟合优度
\[
R^2 = \frac{ESS}{TSS} = 1 - \frac{RSS}{TSS} \in [0,1],
\]
即 $y$ 的总变异中被 $x$ 线性解释的比例：模型完全解释时 $RSS = 0$、$R^2 = 1$，完全未解释时 $R^2 = 0$。\textbf{$t$ 检验。}检验单个系数的显著性需要系数的抽样分布。在正态性与同方差假设下，斜率估计量 $\hat b$ 是 $\varepsilon_i$ 的线性函数，因而服从正态分布
\[
\hat b \sim N\Bigl(b,\ \frac{\sigma^2}{\sum_i (x_i - \bar x)^2}\Bigr),
\]
其方差与扰动方差 $\sigma^2$ 成正比、与 $x$ 的变异成反比：$x$ 的取值越分散，拟合直线的斜率定位越精确。由于 $\sigma^2$ 未知，以 $s^2 = RSS/(n-2)$ 估计，得到标准误 $se(\hat b) = s / \sqrt{\sum_i (x_i - \bar x)^2}$。在 $H_0:\ b = 0$ 下，标准化后的统计量
\[
t = \frac{\hat b}{se(\hat b)} \sim t(n-2),
\]
即估计量偏离零的距离以标准误为单位度量。样本量较大时 $t$ 分布近似标准正态：\important{5\% 显著性水平下 $|t| > 1.96$ 拒绝原假设}，系数统计显著异于零。多元回归中自由度相应改为 $n$ 减去待估参数的个数，检验逻辑不变。
\end{derivation}

$R^2$ 与 $t$ 检验各有明确的边界。$R^2$ 的陷阱：增加任何解释变量都使 $R^2$ 不减（$RSS$ 只会下降），因此多元回归常报告调整 $R^2$；\important{$R^2$ 度量拟合而非因果，高 $R^2$ 不构成任何因果结论的证据}，时间序列中尤须警惕。\mn{记忆口诀：$R^2$ 看"解释了多少"，$t$ 看"是否非零"，两者都不能证明因果；因果由研究设计保证，不由回归输出保证。}$t$ 检验的边界：$|t| > 1.96$ 只说明估计精度足以区分零，不说明效应大小重要（经济显著性与统计显著性是两回事），多重假设检验下还应调整显著性标准。

\subsection{内生性与遗漏变量偏误}

OLS 的无偏性依赖 $\mathbb{E}[\varepsilon_i \mid x_i] = 0$，这一假设在数字经济研究中尤其脆弱：解释变量多为政策、渗透率、指数等选择行为的结果，与扰动项相关的情形大量存在。解释变量与扰动项相关的现象称为内生性，本节梳理其三个来源，并推导遗漏变量偏误的公式。

\begin{definition}[内生性（Endogeneity）]
内生性是指解释变量与扰动项相关的现象，形式上表现为 $Cov(x, \varepsilon) \neq 0$。此时 OLS 估计量不再收敛于真实参数，估计结果不能作因果解释。
\end{definition}

内生性有三个经典来源，每个来源在数字经济研究中都有对应情形。\textbf{遗漏变量}：与解释变量相关的因素被留在扰动项中。估计数字化程度对生产率的影响时，企业管理能力既提高生产率又促进数字化，若不可观测，$Cov(x, \varepsilon) > 0$。\textbf{测量误差}：解释变量的度量含噪声。数字化的度量（是否有网站、线上销售占比）只是真实数字化程度的噪声代理，度量误差进入扰动项。\textbf{反向因果}：被解释变量同时影响解释变量。生产率高的企业更有余力投入数字化，数字支付的发展既促进电商、也被电商需求拉动。三个来源的方向不同，需要分别处理。

\begin{derivation}{遗漏变量偏误公式}{}
\textbf{遗漏变量偏误的推导。}设真实模型为
\[
y = \beta x + \gamma z + u, \qquad \mathbb{E}[u \mid x, z] = 0,
\]
其中 $z$ 为遗漏变量，$\gamma$ 为其对 $y$ 的真实影响，$u$ 为真正的随机扰动（零条件均值意味着 $Cov(x,u) = 0$）。估计者不知道 $z$ 的存在（或没有 $z$ 的数据），错误地只做短回归 $y = \beta x + \varepsilon$。此时 $z$ 的效应被并入扰动项：$\varepsilon = \gamma z + u$，而 $z$ 与 $x$ 相关，故 $Cov(x, \varepsilon) = \gamma\, Cov(x, z) \neq 0$，零条件均值假设失效。\textbf{第一步，短回归系数公式。}一元回归的 OLS 系数是 $x$ 与 $y$ 的样本协方差除以 $x$ 的样本方差：
\[
\hat\beta_{\text{短}} = \frac{Cov(x, y)}{Var(x)}.
\]
\textbf{第二步，代入真实模型。}把 $y = \beta x + \gamma z + u$ 代入分子，利用协方差的线性性逐项展开：
\[
\hat\beta_{\text{短}}
= \frac{Cov(x,\, \beta x + \gamma z + u)}{Var(x)}
= \beta + \gamma \cdot \frac{Cov(x, z)}{Var(x)} + \frac{Cov(x, u)}{Var(x)}.
\]
\textbf{第三步，化简。}$\mathbb{E}[u \mid x,z] = 0$ 保证 $Cov(x,u) = 0$，最后一项消失，得到
\[
\hat\beta_{\text{短}} = \beta + \gamma \cdot \frac{Cov(x, z)}{Var(x)}.
\]
第一项是 $x$ 的真实效应 $\beta$；第二项就是\important{遗漏变量偏误 $\gamma \cdot Cov(x,z)/Var(x)$}：它不随样本增大而消失，数据再多短回归也不会趋近真实参数。其符号由两项共同决定，$\gamma$（$z$ 对 $y$ 的真实影响方向）与 $Cov(x,z)$（$z$ 与 $x$ 的相关方向）；其大小还取决于 $z$ 与 $x$ 相关的强度。
\end{derivation}

偏误公式给出了实证论文最常用的技能：方向性判断。\important{当 $\gamma > 0$ 且 $Cov(x,z) > 0$ 时 OLS 高估，$\gamma > 0$ 且 $Cov(x,z) < 0$ 时低估}。数字经济研究的典型应用：估计"数字化程度对生产率的影响"时，企业能力是典型的遗漏变量，能力强（$\gamma > 0$）的企业更可能数字化（$Cov(x,z) > 0$），简单 OLS 系统性高估数字化的生产率效应；\mn{偏误方向口诀：$\gamma$ 与 $Cov(x,z)$ 同号则高估、异号则低估；答"遗漏变量如何影响估计"先定两个符号，再下方向结论。}教育回报的经典例同理，能力被遗漏使回报率被高估。补救手段分两级：加入控制变量、控制固定效应只能处理可观测或时间不变的遗漏，不可观测且时变的遗漏必须依靠 13.2 节的识别策略。

测量误差的偏误方向与遗漏变量相反。若解释变量的度量含噪声（数字化的度量天然如此：以"企业是否注册域名""线上销售额占比"等指标代理数字化程度，每个代理都只捕捉真实数字化的一部分），噪声会把估计系数向零拉（\important{衰减偏误}），数字化效应被系统性低估；这也是实证报告中"不显著"的一种可能来源，有时是度量问题而非效应不存在。\mn{衰减偏误直观理解：解释变量"测不准"，系数被拉向零；实证结果不显著时，先排查度量精度，再怀疑效应本身。}被解释变量一侧则相反：若 $y$ 的度量误差与 $x$ 无关，它只增大估计方差、不造成偏误，偏误只来自解释变量一侧的噪声。

反向因果不需要推导偏误公式，其逻辑一目了然：$y$ 影响 $x$ 时，回归把双向关系压成单向系数。数字经济中最典型的反向因果是规模与数字化的互促：平台规模扩大带来更多广告与数据资源，又反过来支撑更大的规模，任何"规模对绩效"的 OLS 都混合了两个方向。反向因果的经典解法是寻找只影响 $x$、不直接受 $y$ 影响的第三变量，这正是 13.2.3 节工具变量的思想。

\begin{misconception}
大数据时代的一个流行误区是"数据量足够大，相关关系可以取代因果推断"。相关是样本协方差的性质，因果是反事实的比较：冰淇淋销量与溺水人数高度相关，但共同的驱动是气温。\textbf{数据量不改变推断的性质，只提高估计精度}；因果结论永远需要研究设计（随机化或准实验），这是本章 13.2 节的中心任务。
\end{misconception}

\section{因果推断识别策略}

内生性无法靠"多加变量"根治，因为不可观测的遗漏永远可以存在。因果推断的出路是研究设计：利用制度背景中的外生变异，构造可比的对照。DID、IV、RDD 分别利用政策的时空变异、外生冲击与阈值分配，先固定三者共享的反事实语言。

\begin{definition}[反事实（Counterfactual）]
反事实是指个体在未接受处理时本会出现的潜在结果。个体 $i$ 的处理效应为 $\tau_i = y_i(1) - y_i(0)$，其中 $y_i(1)$ 为接受处理的结果，$y_i(0)$ 为未接受处理的反事实结果；同一时期二者只能观测其一，因果推断的困难正在于反事实不可观测。
\end{definition}

基于反事实语言，可以区分三个效应概念。\textbf{平均处理效应}（Average Treatment Effect, ATE）：全体个体的平均处理效应；\textbf{处理组平均处理效应}（ATT）：实际接受处理者的平均效应；\textbf{局部平均处理效应}（LATE）：只对特定边际群体成立的效应（13.2.3 节）。每个识别策略只能识别其中一个，且要求各自不同的假设，这构成本节的主线。

\subsection{双重差分（DID）}

政策评估的经典问题是：一项数字经济政策（试点、补贴、牌照）推出后，如何把政策效应从时间趋势中分离出来？处理组的变化包含"政策效应 + 趋势"，对照组的变化只有"趋势"，两次做差即可抵消趋势，这正是双重差分的思想。DID 适用于政策只在部分对象、部分时间实施的场合，是数字经济政策评估中使用频率最高的识别策略。

\begin{definition}[双重差分（Difference-in-Differences, DID）]
双重差分是利用政策在"组别 × 时间"两个维度的变异识别处理效应的方法：以处理组与对照组在政策前后的各自变化做差，在\textbf{平行趋势假设}（Parallel Trends Assumption）下抵消共同的时间趋势，识别政策的处理效应。
\end{definition}

\begin{derivation}{DID 估计量}{}
\textbf{两期两组设定下的 DID。}设政策在 $t^*$ 时刻实施，处理组 $T$（受政策影响）与对照组 $C$（不受影响），结果变量在政策前（$t = 0$）与政策后（$t = 1$）各观测一期。\textbf{第一步，处理组的单差。}处理组政策前后的变化为
\[
\Delta_T = \mathbb{E}[y_{T,1}] - \mathbb{E}[y_{T,0}],
\]
它包含两部分：政策本身的效应，以及所有个体共同经历的时间趋势（宏观环境、季节性、技术外溢等与政策无关的变化），两者无法在单差中分开。\textbf{第二步，对照组的单差。}对照组同期变化为
\[
\Delta_C = \mathbb{E}[y_{C,1}] - \mathbb{E}[y_{C,0}],
\]
对照组未受政策影响，其变化仅包含时间趋势，恰好提供了趋势的估计。\textbf{第三步，双重差分。}两次做差：
\[
ATT = \Delta_T - \Delta_C
= \bigl(\bar y_{T,1} - \bar y_{T,0}\bigr) - \bigl(\bar y_{C,1} - \bar y_{C,0}\bigr),
\]
处理组单差减去对照组单差，时间趋势被抵消，剩下的即政策效应。抵消成立的条件是\textbf{平行趋势假设}：若政策不发生，处理组将沿对照组的相同趋势变化，
\[
\mathbb{E}[y_{T,1} - y_{T,0} \mid \text{无政策}] = \mathbb{E}[y_{C,1} - y_{C,0}],
\]
也就是说，对照组的变化就是处理组反事实变化的无偏估计。\important{平行趋势下，双重差分估计量识别政策的 ATT}。
\end{derivation}

DID 识别的不是处理组的绝对变化，而是处理组相对对照组的\important{相对变化}，任何与政策同期、且只冲击处理组的其他因素都会混入估计。因此 DID 对政策时点的外生性要求很高：政策不能是处理组"变差之后"才被动推出的。

DID 在数字经济政策评估中的应用极为丰富。\textbf{"宽带中国"示范城市}：2014 年起工业和信息化部会同国家发展改革委分批确定示范城市（城市群），以入选城市为处理组、未入选城市为对照组，比较电商发展、数字创业等结果变量的政策前后差异，是数字基础设施效应的标准 DID 设计。\textbf{数字人民币试点}：试点城市的分批扩容（第 10 章）提供了天然的时空变异，文献以其估计数字人民币对支付习惯与金融包容的影响。\mn{DID 口诀：两次做差、趋势抵消；事前平行、事后可信。平行趋势用 13.2.2 节的事件研究图检验，政策前系数不显著是 DID 成立的前提。}\textbf{电子商务进农村综合示范}：2014 年起分批实施的示范县政策，常被用于识别电商对农村收入与创业的因果效应（第 9 章三层鸿沟的实证对应）。三类设计的共同软肋也一致：试点批次的选择常与地方禀赋相关（有条件的地区先入选），平行趋势必须事前检验，而不能假定。

\subsection{事件研究法与平行趋势检验}

平行趋势假设本身无法直接检验：反事实趋势不可观测。标准的近似检验是\textbf{事件研究法}（Event Study）：以政策前一期为基期，把政策前后各期的处理组与对照组的差异逐一估计出来，观察这些差异是否在政策前已经出现。事件研究图同时承担另一项任务：刻画政策效应的动态路径，回答"效应是即期的还是渐进的、是否存在预期效应"。

图~\ref{fig:eventstudy} 给出了事件研究图的标准形态：横轴为相对政策时点 $k$（$k < 0$ 政策前，$k \geq 0$ 政策后），纵轴为各期处理组与对照组的差异估计及其置信区间，政策时点处画竖线分隔。读图规则：\textbf{政策前}（$k < 0$）系数应围绕零波动，若出现趋势性差异或个别系数显著，平行趋势假设存疑；\textbf{政策后}（$k \geq 0$）系数跳升表明政策即时生效，逐期上升表明效应渐进释放；\mn{事件研究图读图口诀：事前平、事后升、盯负一（$k=-1$）。事前平是平行趋势的证据，$k=-1$ 显著是预期效应。}\textbf{预期效应}：$k = -1$ 的系数显著提示政策公布先于实施（企业提前调整行为），文献常将基期前移或剔除存在预期效应的样本。

\begin{figure}[htbp]
\centering
\begin{tikzpicture}[>=Stealth]
\begin{axis}[
    width=11cm, height=6cm,
    xlabel={相对政策时点 $k$}, ylabel={各期处理效应估计},
    xmin=-4.8, xmax=4.8, ymin=-0.4, ymax=1.1,
    xtick={-4,-3,-2,-1,0,1,2,3,4},
    x tick label style={font=\small},
    y tick label style={font=\small},
    x label style={font=\small}, y label style={font=\small},
    axis lines=left,
]
\draw[gray, thin] (axis cs:-4.8,0) -- (axis cs:4.8,0);
\draw[dashed, gray, thick] (axis cs:-0.5,-0.4) -- (axis cs:-0.5,1.1);
\node at (axis cs:-0.5,0.92) [rotate=90, font=\small] {政策时点};
\addplot+[color=main, only marks, mark=*, mark size=1.4pt,
    error bars/.cd, y dir=both, y explicit, error bar style={color=main, thin}]
  coordinates {
    (-4, 0.06) +- (0, 0.18)
    (-3,-0.10) +- (0, 0.18)
    (-2, 0.04) +- (0, 0.18)
    (-1, 0.00) +- (0, 0.18)
    ( 0, 0.55) +- (0, 0.22)
    ( 1, 0.75) +- (0, 0.22)
    ( 2, 0.86) +- (0, 0.22)
    ( 3, 0.82) +- (0, 0.22)
    ( 4, 0.79) +- (0, 0.22)
  };
\end{axis}
\end{tikzpicture}
\caption{事件研究图：政策前系数接近零，政策后跳升}
\label{fig:eventstudy}
\end{figure}

事件研究已成为数字经济政策论文的标配报告。中国文献中的标准做法：以政策实施前一期为基期，报告各期系数与 95\% 置信区间，事前系数不显著且无趋势作为平行趋势的证据；数字人民币试点、宽带试点等分批扩容政策的研究均按此程序报告动态效应（第 10 章的制度背景由此进入实证）。平行趋势假设的检验由此有了标准做法：报告事件研究图，以事前系数不显著且无趋势作为支持性证据。

\subsection{工具变量（IV）}

当解释变量内生且控制变量无法根治时，需要寻找\textbf{工具变量}：一个与内生解释变量相关、但只通过解释变量影响被解释变量的外部变量。IV 的思想是"清洗"：只保留 $x$ 中由工具变量驱动的那部分变异，用这部分外生变异估计因果效应。

\begin{definition}[工具变量（Instrumental Variable, IV）]
工具变量是满足两个条件的变量 $z$：\textbf{相关性}，$Cov(z, x) \neq 0$，与内生解释变量相关；\textbf{外生性}（排除性约束），$Cov(z, \varepsilon) = 0$，只通过解释变量影响被解释变量。
\end{definition}

\begin{derivation}{两阶段最小二乘}{}
\textbf{两阶段最小二乘（2SLS）的推导。}设模型 $y = \beta x + \varepsilon$，$x$ 内生（$Cov(x,\varepsilon) \neq 0$），工具 $z$ 满足相关性与外生性。2SLS 的思想是把 $x$ 拆成两半：由 $z$ 决定的部分，与造成内生性的残余部分，只用前者估计 $\beta$。\textbf{第一阶段。}以内生变量对工具变量回归，一元回归的拟合值为
\[
\hat x = \hat\pi\, z, \qquad \hat\pi = \frac{Cov(z, x)}{Var(z)},
\]
$\hat x$ 是 $z$ 的线性函数，而 $z$ 与 $\varepsilon$ 不相关，因此 $\hat x$ 与 $\varepsilon$ 不相关，它是 $x$ 中"干净"的那一半。\textbf{第二阶段。}以 $\hat x$ 替代 $x$ 做回归：
\[
\hat\beta_{IV} = \frac{Cov(\hat x, y)}{Var(\hat x)}.
\]
\textbf{化简。}把 $\hat x = \hat\pi z$ 代入，分子分母的 $\hat\pi$ 消去：
\[
\hat\beta_{IV} = \frac{\hat\pi\, Cov(z, y)}{\hat\pi^2\, Var(z)} = \frac{Cov(z, y)}{Cov(z, x)},
\]
直观含义：IV 系数等于 $z$ 与 $y$ 的协方差除以 $z$ 与 $x$ 的协方差，即"$z$ 每改变 $x$ 一个单位，$y$ 随之改变多少"。\textbf{验证。}代入 $y = \beta x + \varepsilon$ 展开分子：
\[
Cov(z, y) = \beta\, Cov(z, x) + Cov(z, \varepsilon),
\]
外生性（$Cov(z,\varepsilon) = 0$）使第二项消失，$\hat\beta_{IV} = \beta$。\important{IV 只用 $z$ 驱动的外生变异识别 $\beta$}。其识别对象是\textbf{局部平均处理效应（LATE）}：只对"被 $z$ 改变处理状态"的边际群体成立，而非全体样本的 ATE。
\end{derivation}

工具变量改变的是部分个体的处理状态：以距离为例，靠近杭州的地区数字金融渗透更快，但"靠近"这一因素只对受距离约束的个体起作用，IV 估计的正是\important{这些边际个体}的效应。因此 IV 结论的推广要谨慎：它是特定工具变量作用下的局部效应，更换工具可能得到不同的 LATE，这不是估计失败，而是识别对象的差异。\important{IV 识别 LATE 而非 ATE，是计量辨析题的常客}。

\begin{definition}[局部平均处理效应（Local Average Treatment Effect, LATE）]
局部平均处理效应是指对受工具变量影响而改变处理状态的边际群体成立的因果效应：仅覆盖这部分个体，不覆盖全体样本。
\end{definition}

工具变量的两个条件中，相关性可检验，外生性不可检验。\textbf{相关性检验}：第一阶段回归的 $F$ 统计量，\important{经验法则为 $F > 10$}，$F$ 过小即为弱工具变量：弱工具下估计量的方差膨胀、置信区间失真，且小样本偏误向 OLS 回归，工具变量的优点随之消失。\textbf{外生性论证}：$Cov(z, \varepsilon) = 0$ 无法用数据直接检验，只能靠经济逻辑与研究设计论证；当工具数量超过内生变量时，可用\textbf{过度识别检验}（J 检验）作有限度的检验，但排除性约束的最终依据仍是故事本身。\mn{IV 口诀：相关要检验（$F > 10$），外生靠论证；识别的是 LATE 不是 ATE，复试答"IV 估计的是谁的效应"时说"服从者"。}数字经济研究的经典工具变量分三类。

\textbf{地理距离类}：张勋等（2019）以各地区到杭州的球面距离作为数字金融发展的工具变量，外生性来自数字金融平台总部区位的特殊性，相关性来自数字金融服务的辐射半径。\textbf{地形类}：地形起伏度、坡度决定网络基础设施的铺设成本，从而影响数字接入，但地形本身不直接影响多数经济结果，常被用作网络覆盖的工具。\textbf{历史类}：历史上的邮电基础设施沿路径依赖塑造当代数字技术扩散，如黄群慧等（2019）以历史邮电数据作为互联网发展工具的研究。三类工具的共性：借助自然或历史因素构造外生变异，以研究设计的说服力弥补不可检验性。

\subsection{断点回归（RDD）}

当政策按某一连续变量的阈值分配时（评分达标即入选、金额达标即减免），阈值两侧的个体在分配变量上几乎相同，处理分配近似随机。断点回归利用这一不连续性：比较恰好位于阈值两侧的个体，估计政策在阈值处的局部因果效应。

\begin{definition}[断点回归（Regression Discontinuity Design, RDD）]
断点回归是利用处理分配在某一连续变量阈值处的不连续识别因果效应的方法：当 $D_i = \mathbb{1}[r_i \geq c]$ 时，比较分配变量在阈值 $c$ 两侧的个体，以结果变量在断点处的跳跃估计处理效应。
\end{definition}

\begin{derivation}{精确断点的局部平均处理效应}{}
\textbf{精确断点（Sharp RDD）的识别。}设分配变量为 $r$，阈值为 $c$，处理状态由规则 $D = \mathbb{1}[r \geq c]$ 完全决定（达标即处理，不达标即不处理，无例外）。潜在结果 $y(1)$、$y(0)$ 满足\textbf{连续性假设}：$\mathbb{E}[y(1) \mid r]$ 与 $\mathbb{E}[y(0) \mid r]$ 在 $r = c$ 处连续，即除政策外没有任何因素在断点处跳跃。\textbf{第一步，写出观测结果。}实际观测到的结果取决于处理状态：$r \geq c$ 的个体观测到 $y(1)$，$r < c$ 的个体观测到 $y(0)$，因此
\[
\mathbb{E}[y \mid r] = \mathbb{E}[y(0) \mid r] + D(r)\cdot\bigl(\mathbb{E}[y(1) \mid r] - \mathbb{E}[y(0) \mid r]\bigr),
\]
第二项只在阈值右侧（$D = 1$）进入，左侧为零。\textbf{第二步，取左右极限。}由于潜在结果的条件期望在断点处连续，右侧极限不受 $D$ 的跳跃干扰：
\[
\lim_{r \downarrow c} \mathbb{E}[y \mid r] = \mathbb{E}[y(1) \mid c],
\qquad
\lim_{r \uparrow c} \mathbb{E}[y \mid r] = \mathbb{E}[y(0) \mid c].
\]
\textbf{第三步，做差。}两式相减，连续性的部分在极限处重合，剩下的就是断点两侧唯一的差异来源（政策）：
\[
\tau_{RDD} = \lim_{r \downarrow c} \mathbb{E}[y \mid r] - \lim_{r \uparrow c} \mathbb{E}[y \mid r] = \mathbb{E}[y(1) \mid c] - \mathbb{E}[y(0) \mid c].
\]
\important{RDD 识别的是结果变量在断点处的跳跃}，即断点处个体的局部平均处理效应：它是 $r = c$ 处（而非全体样本）的因果效应，因为极限运算只利用断点附近的信息。
\end{derivation}

RDD 的识别假设有清晰的检验路径。\textbf{连续性假设}是核心：任何其他随 $r$ 连续变化的因素都会被极限差消去，只有与政策同时跳变的因素才构成威胁，因此检验的重心是"断点处是否还有其他跳变"。\textbf{操纵检验}：若个体能操纵分配变量使自己恰好落在阈值一侧（如刷分过线），断点两侧的个体密度会出现不连续，文献有标准的密度检验检查断点两侧个体密度是否连续。\textbf{带宽选择}：极限只在 $c$ 处成立，估计需限定在 $c$ 附近的\textbf{带宽}（bandwidth）内，带宽越小偏误越小、方差越大。若断点处处理概率只有跳跃而非从 0 到 1（\textbf{模糊断点}，fuzzy RDD），识别量还需除以处理概率的跳跃。

RDD 的应用场景广泛。中国语境下的经典样本是\textbf{高考录取线}：一分之差决定录取与否，录取线两侧的考生在能力上几乎无差，文献以此研究教育回报，恰好构成 RDD 所需的近似随机分配。\mn{RDD 案例锚点：高考录取线，一分之差决定录取、两侧能力几乎无差，这是"阈值近似随机"的中国经典样本；复试举例先报它。}数字经济中的 RDD 场景同样丰富：信贷平台的\textbf{信用分门槛}（评分达到阈值才获得授信，比较门槛两侧借款人的违约率识别信贷效应）、电商的\textbf{满减门槛}（满减金额附近的订单处理概率发生跳跃）、平台补贴的\textbf{收入门槛}（资格线两侧的用户行为差异）。RDD 的局限也须讲明：它识别的是断点附近的局部效应，远离断点的个体与断点处个体差异悬殊，\important{断点外推广依赖外推，超出 RDD 本身的识别范围}。

\begin{misconception}
两个常见误读。其一，把 RDD 的"局部平均处理效应"解读为全局平均效应：RDD 只识别断点附近，政策对全体对象的平均效应无法由 RDD 单独回答。其二，把阈值两侧的个体视为"完全随机"：阈值分配只是近似随机，连续性假设仍需论证与检验（密度检验），不是自动成立。
\end{misconception}

\subsection{匹配与合成控制}

DID 需要平行趋势与对照组，但两类问题超出其能力：可观测特征上的选择差异，以及只有一个处理对象、找不到对照组的政策评估。前者由倾向得分匹配处理，后者由合成控制法处理。

\begin{definition}[倾向得分匹配（Propensity Score Matching, PSM）]
倾向得分匹配是以可观测特征估计个体进入处理组的概率（倾向得分 $e(x) = P(D = 1 \mid x)$），并为每个处理个体匹配特征相近的对照个体的方法，使两组在可观测维度上可比。
\end{definition}

PSM 的思想：以倾向得分为标尺，为每个处理个体匹配得分相近的对照个体，把按多个特征匹配的问题降为按一个标量匹配的问题。其局限同样清晰：PSM 只能平衡\important{可观测}特征，若选择进入处理组的依据包含不可观测因素（能力、偏好），匹配后的比较仍然有偏；匹配质量还要求对照组中存在特征相近的个体，找不到相近者时只能剔除。因此 PSM 常与 DID 结合使用（PSM--DID）：先匹配再做差，用匹配处理可观测选择、用差分处理时间趋势。当政策只作用于一个对象时，连匹配对象都找不到，合成控制法为这一场景设计。

\begin{definition}[合成控制法（Synthetic Control Method, SCM）]
合成控制法是以对照组个体的加权组合构造"合成的反事实"的方法：选取权重 $w_j \geq 0$（$\sum_j w_j = 1$）使政策前合成单元的结果路径与处理单元精确拟合，比较政策后处理单元与合成单元的差异。
\end{definition}

合成控制法专门服务"只有一个处理对象"的评估：一个省的试点、一个平台的改革都没有天然对照组，但可以由多个未处理对象的加权组合构造一个虚拟对照。经典的示范应用是 Abadie 等（2010）对加州控烟政策的评估：以其他州加权构造"合成加州"，政策前吸烟率路径与加州几乎重合，政策后的偏离即政策效应。数字经济的对应场景：某城市率先实施平台经济试点或数据交易试点，可用其他城市构造"合成城市"评估试点效应。三种方法的适用分工取决于数据的制度背景：\important{识别策略的选择由数据的制度背景决定，而非由方法的难易决定}。有分批试点用 DID，有强外生变量用 IV，有明确阈值规则用 RDD，只有单一对象用合成控制，可观测选择为主用 PSM。

\section{实证案例教学：从问题到结论}

方法的价值在应用。实证研究的完整流程可以概括为四步：\textbf{分析问题}（把经济现象转化为可检验的实证命题）、\textbf{建模}（把命题翻译为模型设定与变量选择）、\textbf{求解}（选择识别策略完成估计与检验）、\textbf{分析}（把估计结果翻译回经济含义）。本节以四个案例完整走一遍这一流程，每个案例都是一份可迁移的研究设计模板：学会的不是案例本身，而是拿到一个新问题时按四步推进的处理路径。

\subsection{案例一：数字化能提高企业生产率吗？}

\textbf{分析问题。}"数字化是否提高企业生产率"是数字经济实证的第一问题（第 8 章增长效应的微观基础）。问题要转化为可检验的命题：给定其他条件不变，数字化程度更高的企业是否拥有更高的全要素生产率？注意"相关"与"提高"的区别：前者只需描述性证据，后者是因果命题，这个区别决定后续每一步的选择。

\textbf{建模。}把命题写成线性模型
\[
y_i = \beta\, x_i + \gamma' C_i + \varepsilon_i,
\]
其中 $y_i$ 为企业全要素生产率的对数，$x_i$ 为数字化程度（如数字技术投入强度），$C_i$ 为控制变量（企业规模、年龄、行业），$\beta$ 是识别的对象：数字化的生产率效应。该设定是 13.1.1 节一元回归的直接推广：解释变量增加为多个，系数的含义不变，仍是"该变量每变动一单位、$y$ 平均变动多少"，其估计同样以残差平方和最小化为准则。

\textbf{求解。}直接用 OLS 的隐患在 13.1.2 节已经预见：能力强（管理、创新）的企业既更数字化、生产率也更高，能力被遗漏在 $\varepsilon_i$ 中，$Cov(x, \varepsilon) > 0$，$\hat\beta$ 系统性高估。修正分两级：加入更多控制变量只能处理可观测部分；根治要换识别策略，以"数据资产入表"（第 7 章）这类外生政策冲击构造 DID，比较入表企业与非入表企业的生产率路径。这一设计的识别难点有二：入表选择内生（率先入表的企业本就数据治理更强），预期效应（入表消息公布时市场已提前反应，入表前的系数可能不为零，事件研究的 $k = -1$ 系数是标准检验点）。

\textbf{分析。}结论表述要区分两种口径："数字化与生产率正相关"是描述性结论，OLS 即可报告；"数字化提高生产率"是因果结论，必须依赖识别策略。若 DID 结果显示入表后生产率显著提升，支持数据要素化的生产率效应；同时要讲清估计的局部性：DID 识别的是"入表"这一政策冲击的效应，不是数字化对生产率的全部效应，外推需谨慎。

\subsection{案例二："宽带中国"试点促进了农村电商吗？}

\textbf{分析问题。}数字鸿沟的因果识别（第 9 章三层鸿沟的实证对应）：网络接入的改善是否促进了农村电商发展？"接入好的地区电商更好"只是相关：禀赋好、动力强的地区可能更早接入，也更容易发展电商，直接比较会混淆两者。

\textbf{建模。}政策冲击："宽带中国"示范城市与电信普遍服务试点自 2014 年起分批推进，先行接入地区为处理组、后行地区为对照组，结果变量取农村电商交易额（或农村人均收入）。政策分批推进提供了组别与时间两个维度的变异，恰好匹配 13.2.1 节的 DID 设定。

\textbf{求解。}按 DID 三步执行：第一步算处理组单差 $\Delta_T$，第二步算对照组单差 $\Delta_C$，第三步两次做差得 $ATT$。用一个示意性算例演示计算逻辑（数字仅为说明方法，非真实数据）：若处理组接入前后电商交易额从 100 升至 130（$\Delta_T = 30$），对照组同期从 90 升至 110（$\Delta_C = 20$），则 $ATT = 30 - 20 = 10$，剔除共同趋势后接入的净效应为 10。求解的关键步骤在检验：以事件研究图（13.2.2 节）检查政策前差异，事前系数不显著才支持平行趋势；试点批次的选择与地方禀赋相关，还需控制地区趋势。

\textbf{分析。}若 $ATT$ 显著为正且平行趋势成立，接入的因果效应成立，为弥合数字鸿沟提供了因果证据而非仅相关性；进一步可报告效应的异质性（不同地区、不同人群的差异），为政策铺开提供依据。识别局限同样要讲清：DID 识别的是"先行接入"的政策效应，试点地区的经验向未试点地区推广时需要额外论证。

\subsection{案例三：数字经济指数怎么构造？}

\textbf{分析问题。}第 1 章的指数测度法回答了"有哪些数字经济指数"，但指数本身怎么造出来？评价数字经济需要把基础设施、产业发展、创新环境等多个维度的指标合成一个总指数，核心难点在赋权：各指标的权重从何而来？主观赋权依赖专家判断、争议大，客观赋权的标准方法是熵值法。

\textbf{建模。}指标构造分三步建模：\textbf{指标选取}（基础设施、产业发展、创新环境等维度选取分项指标）、\textbf{标准化}（各指标量纲不同，先无量纲化，如极差标准化 $x_{ij}' = (x_{ij} - m_j)/(M_j - m_j)$，其中 $m_j$、$M_j$ 为指标 $j$ 的最小值、最大值）、\textbf{赋权合成}（确定各指标权重并加权求和）。熵值法解决的是第三步。

\textbf{求解。}代入数据按以下步骤完成赋权计算。

\begin{derivation}{熵值法赋权}{}
\textbf{熵值法的权重推导。}设 $n$ 个评价对象、$m$ 个指标，$x_{ij} > 0$ 为标准化后的指标值。赋权要回答的问题是：哪些指标携带更多区分对象的信息？熵值法的答案是看指标的离散程度。\textbf{第一步，计算比重。}第 $j$ 个指标下第 $i$ 个对象的比重为
\[
p_{ij} = \frac{x_{ij}}{\sum_{i=1}^{n} x_{ij}},
\]
把该指标在各对象间的取值规范化为一组非负且和为 1 的数，$p_{ij}$ 可视为指标 $j$ 取值在对象间的分布。\textbf{第二步，计算信息熵。}以熵度量该分布的离散程度，指标 $j$ 的信息熵为
\[
e_j = -k \sum_{i=1}^{n} p_{ij} \ln p_{ij}, \qquad k = \frac{1}{\ln n},
\]
常数 $k$ 的作用是把熵规范到 $[0,1]$ 区间，便于跨指标比较。两个极端情形给出熵的直观含义：所有对象取值相同（$p_{ij} = 1/n$，指标完全没有区分力）时，$e_j = -k \cdot n \cdot (1/n)\ln(1/n) = 1$ 取到最大值；取值集中于一个对象（区分力最强）时 $e_j$ 趋于 0。熵越小，指标取值越离散，携带的区分信息越多。\textbf{第三步，确定权重。}以 $1 - e_j$ 度量指标 $j$ 的信息量，归一化后得到权重：
\[
w_j = \frac{1 - e_j}{\sum_{j=1}^{m} (1 - e_j)},
\]
分母是全部指标信息量之和，权重之和为 1。综合得分为各指标标准化值与权重的加权和 $s_i = \sum_j w_j x_{ij}$。\important{熵值法的核心：以变异程度度量信息量，变异越大的指标权重越高}。
\end{derivation}

\textbf{分析。}客观赋权的逻辑：以数据的变异程度为信息量的代理，各地区差异大的指标（如数字基础设施密度）获得高权重，各地区趋同的指标权重低，权重完全由数据决定，避免了主观争议。其局限："信息量"不等于"重要性"，某指标可能变异小但政策意义大（如数据安全指标），熵值法会系统性低估其权重；实践中常将熵值法与专家赋权结合，或对权重施加约束。\mn{熵值法口诀：变异越大、权重越高；"信息量"是统计概念，"重要性"是政策概念，两者不可混淆。}实践上，第 10 章引用的北大数字普惠金融指数即遵循"标准化、赋权、逐级合成"的技术路线，是客观赋权指数的代表性应用；指数质量的检验则回到本章的方法论：指数造好之后，用它估计的政策效应仍需 DID 等识别策略来支撑。

\subsection{案例四：平台的新功能有效吗？}

\textbf{分析问题。}平台经营中的高频问题：新界面、新推荐算法、新补贴方案是否有效？"有效"必须是因果口径：使用新功能的用户表现更好，不等于新功能有效，因为使用者本身就是选择性样本（更活跃的用户更早使用新功能）。对照组必须是"本来会使用、但因随机分组而未使用"的用户，两组差异才是新功能本身的效应。

\textbf{建模。}把用户随机分为实验组与对照组，实验组施加新功能、对照组维持原状，这是\textbf{A/B 实验}的标准设定，其统计逻辑与随机对照试验完全一致。

\begin{definition}[A/B 实验（A/B Testing）]
A/B 实验是将用户随机分为实验组与对照组、对两组施加不同处理（版本 A 与版本 B）并比较结果差异的实验方法，其统计逻辑与随机对照试验（Randomized Controlled Trial, RCT）完全一致。
\end{definition}

随机化保证了两组可比：随机分组使实验组与对照组在可观测与不可观测特征上均衡，任何组间结果差异都不能归因于事前特征，只能归因于处理本身，组间差异即为平均处理效应。这是 13.2 节反事实框架中唯一能直接观测反事实的设计：对照组就是实验组的反事实。平台实验的规模远超传统实验经济学：线下实验室的样本以百人计，平台实验可以以百万用户计，且实验的边际成本几乎为零，因为分组与数据记录都由平台的基础设施自动完成。建模还有两个前提：个体处理值稳定假设（SUTVA），以及随机分组本身的有效执行。

\begin{definition}[个体处理值稳定假设（SUTVA）]
个体处理值稳定假设（Stable Unit Treatment Value Assumption，简称 SUTVA）是指个体结果只取决于自身处理状态、不受他人处理状态影响的假设，是随机化实验统计推断的基础前提。
\end{definition}

\textbf{求解。}比较组间平均差异即可得到平均处理效应；推断的细节在数字经济中有三个特有干扰。\textbf{其一，网络外部性干扰 SUTVA}：平台上的个体处理效应通过社会互动外溢（实验组用户的行为影响对照组用户），SUTVA 被违背，需以聚类（用户群组、城市、商户圈）为单位随机化，并相应调整推断。\textbf{其二，长期效应与短期代理}：实验通常以短期指标（点击率、转化率）为目标，长期效应（用户留存、口碑）需另行评估，短期最优策略可能损害长期价值。\textbf{其三，伦理与知情}：平台实验的对象通常不知情，算法干预（价格、推荐）的实验引发公平性质疑。标志性事件是 2014 年发表于《美国国家科学院院刊》的 Facebook 情绪传染实验：在约 69 万用户（美国国家科学院院刊，2014）不知情的情况下操纵其信息流的情绪基调，结果显示情绪在社交网络中具有传染性，同时引爆了关于平台实验伦理的长期讨论。

\textbf{分析。}实验结论的内部效度高：随机化解决了选择偏差；外部推广仍需讨论，特定用户群、特定时期的实验结果未必适用于其他场景。A/B 实验对经济学的意义超出平台经营本身：平台成为社会科学实验的"实验室"，实验经济学从线下实验室迁入线上平台，\important{因果推断从"利用自然变异"扩展到"制造可控变异"}，这是数字经济对经济学方法论本身最深刻的影响之一。A/B 实验与 RCT 共享随机化逻辑，两者的因果性来自同一来源；网络外部性、长期效应与伦理问题则是数字经济语境下的三个特有约束。

\begin{misconception}
"A/B 实验的结果一定可靠"是常见误区。三个失效通道：SUTVA 违背（用户间相互影响使组间差异偏离真实处理效应）；短期指标与长期价值的背离（点击率上升而留存率下降）；局部结论的外推（特定用户群、特定时期的实验结果未必适用于其他场景）。实验的随机化解决\textbf{内部效度}（internal validity），不自动解决\textbf{外部效度}（external validity）。
\end{misconception}

除上述四个完整案例外，第 6 章理论的实证对应同样可以套用本节四步流程。平台排他性的竞争效应：以某平台引入排他协议为准自然实验，DID 比较协议签订前后商家多归属率与佣金率的变化，难点在论证协议时点的外生性。算法定价与合谋证据：以高频价格数据检测超竞争加成的结构断点，难点在区分"算法合谋"与"平行定价"的合法结果。这两个案例同样可以按四步流程展开：问题定义清楚之后，建模与求解各归其位。

\begin{chapterbox}
\textbf{本章考点清单}

\begin{enumerate}
\item OLS 正规方程的推导（$\hat b = \widehat{Cov}(x,y)/\widehat{Var}(x)$）与无偏性条件 $\mathbb{E}[\varepsilon_i \mid x_i] = 0$ \important{【基础】}
\item $R^2$ 的平方和分解与 $t$ 检验逻辑 \important{【基础】}
\item 内生性三来源与遗漏变量偏误公式 $\gamma \cdot Cov(x,z)/Var(x)$ \important{【基础】}
\item DID 估计量与平行趋势假设 \important{【专硕·扩展】}
\item 事件研究图与平行趋势检验 \important{【专硕·扩展】}
\item IV/2SLS、LATE 与弱工具变量（$F > 10$）\important{【专硕·扩展】}
\item RDD 的断点跳跃、连续性假设与局部效应 \important{【专硕·扩展】}
\item PSM 与合成控制的适用分工 \important{【专硕·扩展】}
\item 案例四步流程：分析问题→建模→求解→分析 \important{【专硕·扩展】}
\item 熵值法赋权：$e_j = -k\sum p_{ij}\ln p_{ij}$ 与 $w_j$ 公式 \important{【专硕·扩展】}
\item A/B 实验的随机化逻辑与 SUTVA \important{【专硕·扩展】}
\end{enumerate}
\end{chapterbox}

方法章的终点正是前沿章的起点：人工智能的经济影响、大模型时代的产业组织、Web3 的治理挑战，这些前沿命题的实证检验都依赖本章的识别策略。第 14 章转入前沿专题，用全书工具分析 AI、大模型与 Web3。
